\documentclass[../../../include/open-logic-section]{subfiles}

\begin{document}

\olfileid{sth}{z}{arbintersections}
\olsection{Apendiks: Tutupan, Komprehensi, lan Irisan}

Ing \olref[sth][z][milestone]{sec}, kita menehi saran
supaya panjenengan nyawang maneh pakaryan naif ing
\olref[sfr][][]{part} lan mriksa yen pakaryan iku bisa
ditindakake ing~$\Zminus$. Yen panjenengan ngetutake
saran iku, ana \emph{siji} bab kang bisa mbingungake:
panggunaan \emph{irisan} nalika Dedekind ngrembug
\emph{tutupan}.

Elinga saka \olref[sfr][infinite][dedekind]{Closure} yen
\begin{align*}
  \closureofunder{f}{o} & = \bigcap\Setabs{X}{o \in X \text{ lan $X$ tutup miturut
$f$}}.
\intertext{Wangun umume yaiku panetepan kaya mangkene:}
  C & = \bigcap\Setabs{X}{\phi(X)}.
\end{align*}
Nanging iki kudu nggawe kita ngati-ati: amarga
Komprehensi Naif ora lumaku, ora ana jaminan yen
$\Setabs{X}{\phi(X)}$ ana. Dadi panetepan kaya mangkene
katon \emph{curang}.

Untunge, panetepan iku ora kudu curang; utawa, yen
wujude saiki pancen \emph{curang}, kita bisa nggarap
kanthi tliti supaya dadi sah. Cara iku wis disinggung ing
\olref[sth][z][sep]{prop:intersectionsexist}, nalika
kita nerangake sebabe $\bigcap A$ ana kanggo saben
$A \neq \emptyset$. Saiki kita bakal ngandharake kanthi
gamblang.

Kanthi Ekstensionalitas, yen kita arep netepake $C$
minangka $\bigcap\Setabs{X}{\phi(X)}$, kang sejatine
kita butuhake mung obyek $C$ kang nyukupi syarat iki:
\begin{equation}\ollabel{bicondelimarbintersection}
	\forall x(x \in C \liff \forall X(\phi(X) \lif x \in X))\tag{*}
\end{equation}
Saiki upamane ana \emph{sawijining} himpunan $S$ kang
nyukupi $\phi(S)$. Supaya entuk
\olref{bicondelimarbintersection}, kita cukup netepake $C$
kanthi \emph{Pamisahan} mangkene:
\[
	C = \Setabs{x \in S}{\forall X(\phi(X) \lif x \in X)}.
\]
Panjenengan bisa mriksa minangka gladhen yen panetepan
iki pancen ngasilake \olref{bicondelimarbintersection}.
% Pilih $x$. Upamane luwih dhisik $x \in C$ miturut panetepan anyar; mula $x
% \in S$ lan $\forall X(\phi(X) \lif x \in X)$; nanging amarga
% $\phi(S)$, iki padha karo syarat $\forall X(\phi(X)
% \lif x \in X)$. Kosok baline, upamane $\forall X(\phi(X) \lif
% x \in X)$; banjur, amarga $\phi(S)$, kita uga ngerti yen $x \in S$; mula
% $x \in C$ miturut panetepan anyar.
Strategi umum iki bisa ngindhari panggunaan
Komprehensi Naif kang katon ana ing panetepan irisan.
Kanggo kasus kang miwiti pamikiran iki, yaiku
$\closureofunder{f}{o}$, carane mangkene. Ing wiwitan
bukti \olref[sfr][infinite][dedekind]{closureproperties},
kita nyathet yen $o \in \ran{f}\cup\{o\}$ lan
$\ran{f} \cup \{o\}$ tutup miturut $f$. Mula, kang
kita karepake bisa ditetepake mangkene:
\[
	\closureofunder{f}{o} = \Setabs{x \in \ran{f} \cup \{o\}}{(\forall X \ni o)(X \text{ tutup miturut $f$} \lif x \in X)}.
\]

\end{document}
